\section{Discussion: Theoretical and Practical Implications}
\label{sec:discussion}

Having presented the four lifecycle audits, we step back and discuss
what the survey adds relative to existing work, what it means for
theory of evaluation in language technology, and what practitioners
can immediately do with it. The discussion is deliberately limited to
the survey's own scope; we do not extrapolate to voice-interaction
design guidance beyond the evaluation, deployment, and governance
questions that the audit addresses.

\subsection{How This Survey Differs From Prior Work}
\label{subsec:discussion_differentiation}

Five recent reviews cover overlapping territory: two Speech-LLM
surveys \citep{cui2024recent,peng2025surveyspeechlargelanguage}, a
dialogue-oriented review \citep{ji2024wavchat}, an earlier speech
foundation-model review \citep{latif2023sparks}, and a
vision-centric large-multimodal-model safety survey
\citep{ye2024survey}. Table~\ref{tab:survey_comparison} summarizes the
comparison. All five treat their topic mainly as construction:
which architectures exist, which training recipes work, which models
achieve which benchmark scores. The present survey inverts that
question. It asks under which conditions the resulting systems can be
trusted, and it forces every section to be read through that lens.
Concretely, the value added by this survey is threefold.

First, the meta-evaluation of Section~\ref{sec:evaluation_benchmarking}
audits 18 benchmarks along four validity axes rather than compiling a
capability catalog. To our knowledge no prior Speech-LLM survey does
this, and the resulting evidence---17 capability dimensions,
$18\times 17$ coverage matrix, four documented validity gaps---is
released as supplementary material so that future benchmark designers
can build on it or contest it.

Second, the deployment synthesis of Section~\ref{sec:deployment}
organizes serving, architecture, streaming, efficiency, and full-duplex
research around five recurring failure modes (latency fragmentation,
streaming mismatch, memory growth, hardware mismatch, interaction
incompleteness). The prior WavChat review \citep{ji2024wavchat} covers
dialogue and duplex behavior in isolation; here they are linked to the
evaluation gaps in Section~\ref{sec:evaluation_benchmarking} and to the
safety surface in Section~\ref{sec:safety_privacy_governance}.

Third, the responsible-design treatment of
Section~\ref{sec:safety_privacy_governance} develops a threat model
that is specific to speech---audio jailbreaks, cloning-based
impersonation, speaker re-identification, and provenance under
transformation---rather than inheriting the vision-centric taxonomy of
prior LVLM safety surveys \citep{ye2024survey}. Watermark robustness in
particular is reframed as a robustness-under-transformation problem
rather than a detection problem, following the SoK conclusion that no
surveyed audio watermarking scheme is robust to all tested distortions
\citep{wen2025sokrobustaudiowatermarking}.

Cutting across all three, the coupling evidence of
Table~\ref{tab:coupling_evidence} makes the lifecycle framing testable
rather than rhetorical: each coupling names an upstream decision, a
downstream stage, and a documented consequence.

\subsection{Theoretical Implications for Language-Technology Evaluation}
\label{subsec:discussion_theoretical}

The main theoretical contribution is to bring \emph{measurement
validity} back into the discussion of Speech-LLM benchmarks. Classical
psychometrics distinguishes construct, ecological, evaluator, and
operational validity as separate threats; language-technology
evaluation has largely collapsed these into a single question of
``does the model reach the state of the art on the benchmark?''. The
audit in Section~\ref{sec:evaluation_benchmarking} shows that this
collapse hides three systematic risks. Construct validity is
threatened by hard-label scoring of intrinsically distributional
targets such as emotion and style. Evaluator validity is threatened
when proprietary judges score the same model family that helped
curate the benchmark, and by only moderate reported human--judge
correlations (Spearman $\rho \approx 0.61$ for SD-Eval; Pearson $\approx
0.8$ for SeaLLMs-Audio). Operational validity is threatened when
latency, streaming state, and hardware are omitted from evaluation
reports even though these determine whether a system is usable in
production.

Bringing these threats back to the surface has two follow-on
implications. Benchmarks and surveys should treat validity as a
first-class object, not as an afterthought in the discussion section.
And, for benchmarks used to select or promote systems, the coupling
between validity and downstream use---who scored what, under which
conditions, and with what latency budget---should be reported
alongside the numerical scores. The five design principles at the end
of Section~\ref{sec:evaluation_benchmarking} operationalize these
implications for future benchmark releases.

\subsection{Practical Implications for Practitioners}
\label{subsec:discussion_practical}

For information-processing researchers and practitioners choosing or
building Speech-LLM systems, the audit yields four concrete pieces of
guidance.

First, benchmark scores alone are insufficient for system selection.
Because the Use-case$\times$Capability cell of
Table~\ref{tab:benchmark_taxonomy} is empty, a practitioner who needs,
for example, an emotion-aware multilingual voice agent cannot rank
candidate models on published leaderboards. The recommended path is
to pair a use-case-grounded benchmark (VoiceBench, VoiceAgentBench,
Audio2Tool, SLUE-PERB) with at least one capability-focused benchmark
covering the paralinguistic or robustness dimension that matters for
the target deployment.

Second, deployment claims should be reported jointly with latency,
streaming policy, and hardware. Section~\ref{sec:deployment} proposes
a Deployment Card checklist covering model size, hardware class,
batch size, streaming policy, input and output duration, active
optimizations, and the interaction regime under which the numbers
were measured. This card is not a new instrument to be built from
scratch; it is a minimum reporting standard that lets readers compare
across papers without re-running each system.

Third, safety evaluation should be treated as multimodal and adaptive.
Section~\ref{sec:safety_privacy_governance} shows that text-only
alignment does not transfer to spoken inputs, that adversarial
waveform perturbations can bypass transcript-level filters, and that
adaptive purification attacks (VocalBridge
\citep{abbasihafshejani2026vocalbridgelatentdiffusionbridgepurification})
can undo perturbation-based privacy defenses. Deployed voice systems
therefore need layered defenses at the input, representation,
transcription, output, and operational levels, and safety benchmarks
should include perturbed and multimodal audio-text attacks alongside
clean prompts.

Fourth, governance documentation should be modality-aware. General
Model Cards \citep{mitchell2019modelcards} and AI Cards
\citep{golpayegani2024aicardsappliedframework} provide useful templates
but do not by themselves cover speech-specific risks. The governance
requirements listed in Section~\ref{sec:safety_privacy_governance}
name five components---data, safety, privacy, provenance, and
operational governance---that a modality-aware speech-system
description should document.

\subsection{Scope of the Discussion}
\label{subsec:discussion_scope}

We deliberately keep the implications narrow. The survey does not
propose a new user-facing design methodology for voice interfaces,
does not measure real-world user preferences, and does not audit
commercial deployment reports beyond what appears in the peer-reviewed
and preprint literature covered above. The reader should treat the
Deployment Card, Evaluation Card, and governance-requirements list as
reviewer-facing framing tools rather than validated measurement
instruments; the empirical validation of these templates against
community use is left to future work
(Section~\ref{sec:future_directions}) and discussed as a limitation in
Section~\ref{sec:limitations}.
